\subsection{The original quantified pinned point-sum conclusion}
\label{sec:archimedean-pinned-row-closure}

The gap threshold in Proposition~\ref{prop:archimedean-ap7-reserve}
is an internal analytic constant. Denote it by \(\Gamma_g\).
It is distinct from the external constant \(C_g\) in the
original arithmetic context. The latter is fixed before the
requested weight, and is not changed in the following argument.
The original quantifiers permit enlargement of the fixed
thresholds \(C'_0,C_{\ell,0}\). This supplies \(\Gamma_g\)
without adding a hypothesis on the row.
\cite{current-construction-record}

\begin{theorem}[The literal pinned generic point sum]
\label{thm:archimedean-pinned-row-closure}
Fix six real constants
\[
                 C_g,c_b,C_w,c_T,C_T,C_0
\]
and a weight \(A\in\mathbb N\). There are constants
\(C_B\ge0,C'_0,C_{\ell,0}\) such that the following holds.
For every \(C'\ge C'_0\), \(C_\ell\ge C_{\ell,0}\), every
original setting \(D,\mathrm{large}_2,X,Y,u\) satisfying
\(\mathrm{Ctx}(C_g,c_b,C_w,c_T,C_T,C_0,C',C_\ell)\),
every original arithmetic mass tie and admissible mass datum,
and every original pinned quotient measure \(\nu\), one has,
for all auxiliary \(\nu_L,n_{\rm gen}\),
\[
 \mathrm{PointSumInput}\!\left(
   \mathrm{genT}(Y.tO(u,u)),\,
   \mathrm{NOfG}(Y,u,\nu_L,\nu,n_{\rm gen}),\,
   A(r+1)^2,\,
   e^{C_Bm^5f\log b}I_{\rm re}\right).
                                                        \tag{PC1}
\]
Here \(b=\mathrm{large}_2.sc.b\), \(m=2r+2\), and
\(I_{\rm re}\) is the real integral in the given admissible
mass datum. Thus (PC1) is the original
\(\mathrm{StepThmBPinnedG}\) statement, with its original
measure and real integral.
\end{theorem}

\begin{proof}
Fix the six external constants and \(A\). Put
\[
 \beta=\max\{2,C_w\},\qquad
 C_h=|C_T|+\tfrac18+2|C_g|,\qquad B=b^\beta.
                                                        \tag{PC2}
\]
Apply Proposition~\ref{prop:archimedean-ap7-reserve} with
requested weight \(a_w=A\) and this \(C_h\).
It gives \(C_{\rm rot}\) and an internal threshold
\(\Gamma_{g,0}\), depending only on the fixed constants
and \(A\). Increase it to
\(\Gamma_g=\max\{\Gamma_{g,0},0\}\), and set
\[
                 C_B=\beta\max\{C_{\rm rot},0\}.
                                                        \tag{PC3}
\]
Choose \(C'_0\) above the already established local-data,
exact-comparison and occupied-point thresholds. Choose
\(C_{\ell,0}\) above their corresponding thresholds and
above \(8\) and \(128\Gamma_g\).
These choices are made before the setting, rank, row, mass
datum, point and measures.

The original size hypotheses imply
\[
       m\ge C_\ell\ge\max\{8,128\Gamma_g\},
                         \qquad k\ge m.
                                                        \tag{PC4}
\]
The literal standard-hypothesis conversion
\(\mathrm{stdOfHyp\_rec}\) requires
\[
\begin{gathered}
 C_\ell\ge8,\quad C_h\ge0,\quad \beta\ge2,
 \quad C_w\le\beta,\quad C_T\le\beta C_h,
 \quad \tfrac18+2C_g\le\beta C_h,\\
              128|C_s|\le m,\qquad |C'_s|\le m.
\end{gathered}                                         \tag{PC5}
\]
All these hold with \(C_s=C'_s=\Gamma_g\), by (PC2)
and (PC4). Its conclusion supplies the original
\(\mathrm{StdHyp}\) at analytic size \(B\), with the
noncompact-gap threshold \(\Gamma_g\), and supplies
\[
       \mathrm{AllGaps}(\lambda_v,\Gamma_g k m^3)
       \quad\Longrightarrow\quad
       \mathrm{AllGaps}(\lambda_v,\Gamma_g m^4)
                                                        \tag{PC6}
\]
at every real place. The implication uses \(k\ge m\).
Moreover \(k\ge8\), so the cloud-gap set required in the
weighted proof is nonempty. Consequently the reserve
and all shifted-family hypotheses of (AR9) follow from
the existing context. In particular no quantified pointwise
statement outside that finite shift family is used.

The actual local data and
Theorem~\ref{thm:flag-exact-norm-comparison} give one
coefficient at each place, the canonical root scale
\(\theta=1\), and one positive constant \(K_0\) such that
\[
 J_{\rm pin}^2(w)
       =K_0^{-1}\prod_v\|\mathcal J_v(w)\|_{\rm HS}^2,
 \qquad
 K_0^{-1}\prod_v jNorm2(F_v)=I_{\rm re}.
                                                        \tag{PC7}
\]
These are exactly the original amplitude and row-norm
comparisons with \(C_a=C_n=C_1=0\), \(K_v=1\), and
\(C_{\rm fac}=K_0^{-1}\). The prescribed unitary phase
rotation identifies their local orbitals with those in
\(BSumR\) and preserves their radial norms.

For a finite set \(S\) of occupied generic row points,
the established row-point dictionary chooses an integral
normal-form point \(w(\gamma)\) above each \(\gamma\in S\).
The choice is injective: equality of normal-form points
implies equality of their original row labels. Every such
point has the required positive Gram at each real place.
The original weight comparison has constant one:
\[
\begin{split}
 &\mathrm{blockMinorWeight}(\gamma)\,n(\gamma)
       \{NG(\gamma)ND(\gamma)\}^{A(r+1)^2}\\
 &\hspace{1cm}\le
 wPt(w(\gamma))\,ptWeightG(w(\gamma))
                  normWt(w(\gamma))^{A(r+1)^2}.
\end{split}                                           \tag{PC8}
\]
It follows from the exact point-density identity and
\(NG\,ND\le normWt\), retaining the original root
discriminant factors. Multiply (PC8) by the nonnegative
pinned orbital square and use (PC7). The finite sum is
bounded by \(K_0^{-1}BSumR\). Applying (AR9) at size \(B\)
and then (PC7) gives
\[
 K_0^{-1}BSumR
 \le B^{C_{\rm rot}m^5f}I_{\rm re}
 \le e^{C_Bm^5f\log b}I_{\rm re}.
                                                        \tag{PC9}
\]
The last inequality uses \(I_{\rm re}\ge0\), which follows
also from (PC7). This argument never divides by that integral,
so the zero-norm case is included.

Finally, in the literal definition of \(\mathrm{NOfG}\),
the fields used in this point sum are
\[
\begin{aligned}
 \mathrm{jacobianSquareBound}&=J_{\rm pin}^2,
       & n&=nPt,\\
 NG&=NGPt, & ND&=NDPt.
\end{aligned}
\]
They do not depend on
\(\nu_L\) or \(n_{\rm gen}\); these auxiliary data enter
other fields. The occupied-point argument and (PC8) hold
for every choice of those auxiliary fields. The coefficients
and \(K_0\) were chosen before them. Thus (PC9) holds with
the full universal quantifiers in (PC1).

There is no circular threshold choice: \(C_h,\beta,C_*\),
\(C_{\rm rot}\) and \(\Gamma_g\) are fixed first from the
six external constants and \(A\); only then are
\(C'_0,C_{\ell,0}\) fixed. The subsequent lower bound on
\(m\) is a consequence of the original context, not an
extra condition imposed after the row is given.
\end{proof}
